% ============================================================
% Appendix A: Timeline of Hilbert's Problems
% ============================================================

\chapter*{Timeline of Hilbert's Problems}
\addcontentsline{toc}{chapter}{Timeline of Hilbert's Problems}

This timeline traces the major events in the history of Hilbert's 23 problems, from their statement in 1900 to the present day.

\section*{Before 1900}

\begin{description}
    \item[1742] Goldbach writes to Euler with his conjecture (Problem 8).
    \item[1788] Lagrange writes \textit{M\'ecanique Analytique}, establishing the calculus of variations (Problems 20, 23).
    \item[1824] Abel proves the impossibility of solving the general quintic by radicals (Problem 13).
    \item[1829] Gauss proves the fundamental theorem of algebra.
    \item[1831] Galois develops the theory of groups and their application to polynomial equations (Problem 13).
    \item[1854] Riemann's Habilitationsschrift on the foundations of geometry (Problems 4, 16).
    \item[1859] Riemann proposes the Riemann Hypothesis.
    \item[1874] Cantor introduces the concept of cardinality and proves the uncountability of the reals (Problem 1).
    \item[1882] Lindemann proves $\pi$ is transcendental (Problem 7).
    \item[1888] Hilbert proves the finite basis theorem for polynomial rings (Problem 14).
\end{description}

\section*{1900: The Problems Are Proposed}

\begin{description}
    \item[August 8, 1900] Hilbert presents his 23 problems at the International Congress of Mathematicians in Paris.
\end{description}

\section*{1901--1920}

\begin{description}
    \item[1901] Dehn solves Problem 3 (the Dehn invariant).
    \item[1904] Bernstein proves regularity for certain elliptic PDEs (Problem 19).
    \item[1904] Fredholm develops the theory of integral equations, leading to Hilbert space.
    \item[1907] Koebe and Poincar\'e independently prove the Uniformization Theorem (Problem 22).
    \item[1909] Hilbert proves the Waring conjecture (related to Problem 8).
    \item[1910--11] Bieberbach proves his theorems on crystallographic groups (Problem 18).
    \item[1910] Formalism: Russell and Whitehead's \textit{Principia Mathematica} (related to Problem 2).
    \item[1913] Hardy and Littlewood develop the circle method (Problem 8).
    \item[1920] Takagi develops class field theory (Problems 9, 12).
\end{description}

\section*{1921--1950}

\begin{description}
    \item[1921] Noether proves the ascending chain condition for polynomial rings (Problem 14).
    \item[1927] Artin proves Hilbert's 17th problem.
    \item[1929] Gelfond proves $e^{\pi\sqrt{163}}$ is transcendental (Problem 7).
    \item[1931] G\"odel proves the incompleteness theorems (Problem 2).
    \item[1932] von Neumann axiomatizes quantum mechanics (Problem 6).
    \item[1933] von Neumann proves continuity implies differentiability for Lie groups (Problem 5).
    \item[1934] Gelfond--Schneider theorem (Problem 7).
    \item[1936] Gentzen proves relative consistency of arithmetic using transfinite induction (Problem 2).
    \item[1937] Vinogradov proves the ternary Goldbach conjecture for large numbers (Problem 8).
    \item[1940] G\"odel constructs the \textit{L} universe, showing CH is consistent with ZFC (Problem 1).
    \item[1941] Pontryagin develops the maximum principle (Problem 23).
    \item[1950] Chen proved his theorem on Goldbach (Problem 8).
\end{description}

\section*{1951--1980}

\begin{description}
    \item[1955] Grothendieck develops sheaf cohomology (Problem 16).
    \item[1956] De Giorgi proves regularity for elliptic PDEs with measurable coefficients (Prob.\ 19).
    \item[1957] Nash proves regularity independently (Problem 19).
    \item[1963] Cohen proves the independence of CH from ZFC using forcing (Problem 1).
    \item[1965] Sydler completes the characterization of scissors congruence (Problem 3).
    \item[1966] Baker proves his theorem on linear forms in logarithms (Problem 7).
    \item[1970] Matiyasevich proves the undecidability of Diophantine equations (Problem 10).
    \item[1973] Chen proves his theorem on Goldbach (Problem 8).
    \item[1974] Deligne proves the Weil conjectures (Problem 16).
\end{description}

\section*{1981--Present}

\begin{description}
    \item[1989] Dranishnikov disproves the Bieberbach conjecture in high dimensions (Problem 18).
    \item[1994] Andrew Wiles proves Fermat's Last Theorem (related to Problem 10).
    \item[1995] The Langlands--Shahidi method advances the Langlands program (Problem 9).
    \item[2000] The Clay Mathematics Institute announces the Millennium Prize Problems.
    \item[2003] Perelman proves the Poincar\'e conjecture (related to Problem 16).
    \item[2013] Helfgott proves the ternary Goldbach conjecture (Problem 8).
    \item[2018] Peter Scholze receives the Fields Medal for work on perfectoid spaces (related to Problems 9, 12).
    \item[2026] The Langlands program continues to be one of the significant areas of mathematical research.
\end{description}
